% Compiler: PDFLatex
% Bibliography: BibTex

\documentclass[
11pt,
footheight=40pt
]{scrartcl}
\usepackage{ngerman}
\usepackage[utf8]{inputenc}
\usepackage[T1]{fontenc}
\usepackage{graphicx}
\usepackage{listings}
\usepackage{lastpage}
\usepackage{float}

\usepackage[]{todonotes} % disable oder leer lassen
\usepackage{dialogue}

\usepackage{multicol}

\usepackage{pdfpages} % mehrseitige PDFs einbinden


\usepackage[dvipsnames]{xcolor}
\usepackage{amsmath}
\usepackage{amssymb}
\usepackage{gensymb}
\usepackage{cancel}

\usepackage{enumitem}
\usepackage{acronym}
\usepackage{listings}
\usepackage{hyperref}
% Hyperref konfigurieren
\hypersetup{
    colorlinks=true,
    linkcolor=DRSblue,
    filecolor=DRSyellow, 
    citecolor=DRSred,     
    urlcolor=DRSblue,
    pdftitle={Overleaf Example},
    pdfpagemode=FullScreen,
}

% Tabellen

\usepackage{setspace} % For row height adjustments
\usepackage{colortbl} % Für Hintergrundfarben in Tabelle
\usepackage{tabularx}
\usepackage{array} % For tables, z. B. für \newcolumntype

% Schriftart
\usepackage{lmodern}
% Für serifenlose Schrift! Für Serifen (Standard) einfach auskommentieren
% https://www.overleaf.com/learn/latex/Font_typefaces
\renewcommand{\familydefault}{\sfdefault}

 % Pfade zu images, listings, boxes
\newcommand{\imgpath}{assets/img}
\newcommand{\filepath}{assets/files}
\newcommand{\lstpath}{assets/listings}
\newcommand{\vorlagenpath}{/Users/florian/workspace/vorlagen/latex/input}
% Hier sind die Farben und Boxen definiert
\input{\vorlagenpath/boxen}
% Hier sind die selbst definierten Listing-Styles und Sprachen definiert
\input{\vorlagenpath/listings}
\input{\vorlagenpath/aufgaben}

% Variablen für Titlepage und Header/Footer
\newcommand{\lernfeldnr}{12}
\newcommand{\lernfeldtext}{Ein Projekt durchführen}
\newcommand{\bildungsgang}{Fachschinformatiker/in Anwendungsentwicklung}
\newcommand{\scripttitel}{Grundkurs UML}
\newcommand{\untertitel}{Begleitskript zum Unterricht}


%%%%%%%%%%%%%%%%%%%%%%%%%%%%%% Page Layout %%%%%%%%%%%%%%%%%%%%%%%%%%%%%%
\usepackage{geometry}
 \geometry{
 a4paper,
 top=35mm,
 left=25mm,
 right=25mm,
 bottom=30mm
 }

 %%%%%%%%%%%%%%%%%%%%%%%%%%%%%% Kopf- und Fußzeilen %%%%%%%%%%%%%%%%%%%%%
\usepackage[headsepline, automark]{scrlayer-scrpage}
\automark{section}
\setheadsepline{1pt}
\pagestyle{scrheadings}

\ohead{\includegraphics[height=1.7cm]{\boxenimages/drslogo}}
\ihead{\scripttitel \\
%\textbf{Begleitskript zur projektorientierten Arbeit} \\
% \bildungsgang \\
 \leftmark}
\chead{}

\ifoot{Autor: Soulier \\ Stand: \today}
\cfoot{Seite \thepage\ von \pageref{LastPage}}
%\ofoot{\thepage}


 %%%%% Definition der Variablen, um Textteile ein- oder auszublenden %%%%%%
% Entscheidung ob mit oder ohne Musterlösung
\newif\ifloesung
%%%%%%%%%%%%%%%%%%%%%%%%%%%%%%%%%%%%%%%%%%%%%%%%%%%%%%%%%%%%%%%%%%%%%%%%%

%%%%%%%%%%%%%%%%%%%%%%%%%%%%%% Einstellungen %%%%%%%%%%%%%%%%%%%%%%%%%%%%
\setlength{\parindent}{0pt}
%\loesungtrue

\newcounter{cntA}  	
\setcounter{cntA}{0}

\begin{document}

%%%%%%%%%%%%%%%%%%%%%%%%%%%%%% Titlepage %%%%%%%%%%%%%%%%%%%%%%%%%%%%%%%
\begin{titlepage}
  \centering
  %\includegraphics[width=4cm]{\boxenimages/drslogo}\par\vspace{2cm}
  %{\scshape\LARGE \lernfeldtext \par}
  \includegraphics[width=14cm]{\imgpath/title}\par\vspace{2cm}
  {\huge\bfseries \scripttitel \par}
  %\vspace{0.5cm}
  {\Large \untertitel \par}
  %\vspace{8cm}
  %{\large \bildungsgang \par}
    \vfill
  {\large Lehrkraft: Soulier\par}

  {Stand: \large\today\par}
\end{titlepage}

\setcounter{page}{1}           % Manuell auf 1 setzen
\pagenumbering{arabic}         % Arabische Ziffern (1, 2, 3, ...)

\begin{importantbox}{Wichtiger Hinweis}
Dieses Skript befindet sich im Aufbau. Es ist weder vollständig, noch ist garantiert, dass es fehlerfrei ist. Informiere dich unbedingt auch im Internet oder in Büchern. Der Lehrer freut sich über Bug Reports!
\end{importantbox}

\begin{importantbox}{Nur für den internen Gebrauch zulässig!}
Dieses Skript ist \textbf{nicht öffentlich}! Bitte nutze das Skript zum Lernen, kritzel hinein, kopiere Teile für deine Notizen usw. Veröffentliche es aber bitte nicht im Internet oder auf den Drive eurer Firma. Im Skript sind einige Quellen noch unbenannt oder einfach nur hinein kopiert und ich muss diese erst durch selbst erzeugte Inhalte ersetzen.
\end{importantbox}

%\begin{notebox}{\dq Lehrplan \dq PC-Führerschein}
%\textbf{Die Schülerinnen und Schüler verfügen über die Kompetenz, ...} \\
% ...
%\end{notebox}

\begin{warningbox}{Nutzung von LLMs (Chat-GPT)}
Die Nutzung von \acp{LLM} können wir als Lehrer weder verbieten noch die Einhaltung von Regeln diesbezüglich kontrollieren. Es ist auch nicht schlimm, dass ihr diese Technologien benutzt. \\

 In eurem eigenen Sinne empfehle ich euch jedoch, die \acp{LLM} sinnvoll und mit Kopf einzusetzen. Nutzt die Tools nicht, um euch das \textbf{Denken abnehmen zu lassen}, sondern um euer \textbf{Denken zu unterstützen}.  \\ 
 
 \textbf{Lernen geschieht durch Selbermachen} und lasst ihr euch jeden kreativen und anstrengenden Prozess abnehmen, werdet ihr deutlich weniger lernen und nachhaltig behalten, als ihr könntet, wenn ihr eure Denkmuskeln benutzt hättet. \\
 
Besteht bei Abgaben der Verdacht, dass die Lösung LLM-generiert ist, behält sich die Lehrkraft eine \textbf{mündliche Prüfung} an der Tafel oder einen Test auf Papier vor.
\end{warningbox}

%\begin{tipbox}{Informationsquellen}
%\begin{itemize}
%	\item Gratzke, Patett, Schäfer \textit{Westermann} \textbf{IT-Berufe Fachstufe II Lernfelder 10-12 für Anwendungsentwicklung und Daten- und Prozessanalyse} \\
%	Hier bekommt man einen fachlichen Überblick
%	\item weitere Informationsquellen folgen!
%	\item \textbf{alte IHK-Prüfungen} \\
%	Es ist essenziell, vor der eigenen Prüfung mal ein paar Abschlussprüfungen gesehen und gelöst zu haben
%\end{itemize}
%\end{tipbox}

\section*{Versionsverzeichnis}
\begin{itemize}
	\item Version 1 (SJ 26/27): im Aufbau, Aktivitäts-, Anwendungsfall-, Sequenzdiagramme
\end{itemize}

\listoftodos
\newpage


\tableofcontents
\newpage
\section*{Abkürzungsverzeichnis}
\begin{acronym}[Bash]
	\acro{OOA}{Objektorientierte Analyse}
	\acro{OOD}{Objektorientiertes Design}
	\acro{OOP}{Objektorientierte Programmierung}
	\acro{LLM}{Large Language Model}
	\acro{UML}{Unified Modelling Language}
\end{acronym}
%\newpage
\listoffigures
%\newpage
\listoftables
%\newpage
\lstlistoflistings
\newpage

\section{Modellierungssprachen}
In der Informatik nutzen wir Modellierungssprachen um die Anforderungen an ein Softwaresystem, als auch dessen innere Abläufe und Strukturen in einer höheren Abstraktionsebene festzulegen. Modellierungssprachen finden Anwendung in der
\begin{itemize}
	\item \textbf{Entwurfsphase}, wenn die Anforderungen ermittelt und festgelegt werden
	\item \textbf{Implementierungsphase}, wenn die Softwareentwickler/innen die Software in einer höheren Programmiersprache umsetzen
	\item \textbf{Dokumentation} für Kunden oder Kollegen (weil sie ja eh schon angefertigt wurden, kann man diese hier prima weiter verwenden)
\end{itemize}
% ------------------------------------------------------------------------------------------
\subsection{Unified Modeling Language}
Die \ac{UML} ist eine grafische Modellierungssprache zur Spezifikation, Konstruktion, Dokumentation und Visualisierung von Software-Teilen und anderen Systemen. Sie ist derzeit in der Major-Version 2 in der ISO 19505 spezifiziert. In Abb. \ref{fig:uml_diagrammhierarchie} ist eine Übersicht über die gängigsten UML-Diagramme dargestellt. UML gliedert sich im Wesentlichen in 
\begin{itemize}
	\item \textit{Strukturdiagramme}, um die Struktur (z. B. Klassen, Komponenten) einer Software zu beschreiben
	\item \textit{Verhaltensdiagramme}, um das Verhalten einer Software (z. B. Aktivitäten oder Zustände) zu beschreiben
	\item \textit{Interaktionsdiagramme}, um die Interaktion zwischen Software-Komponenten, aber auch zwischen Nutzern zu modellieren
\end{itemize}

\begin{figure}[H]
	\centering
	\includegraphics[width=1\textwidth]{\imgpath/uml_diagrammhierarchie.png}
	\caption{UML-Diagrammhierarchie (Wikipedia)}
	\label{fig:uml_diagrammhierarchie}
\end{figure}

% ------------------------------------------------------------------------------------------

\subsection{Draw.io einrichten}

\begin{tipbox}{Software}
	\textbf{Bevor du beginnst:} Nutze zur Modellierung die Online-Software \textbf{draw.io}. Blende zuvor die UML 2.5 Symbole ein und stelle sicher, dass du im Diagramm die \textbf{richtigen Symbole} verwendest!	
\end{tipbox}

\begin{importantbox}{Symbole}
	Wir sind Profis. Die Nutzung der richtigen Symbole gehört genauso zu unserer Professionalität wie die fachliche/inhaltliche Richtigkeit unserer Arbeit.	
\end{importantbox}

\begin{figure}[H]
	\centering
	\includegraphics[width=1\textwidth]{\imgpath/drawio_shapes.jpg}
	\caption{Draw.io UML 2.5 Symbole aktivieren}
	\label{fig:drawio_shapes}
\end{figure}

\newpage
\section{Anwendungsfalldiagramm}

\begin{notebox}{Einsatzzweck}
	Anwendungsfalldiagramme (Use Case Diagrams) erlauben die Beschreibung von Einsatzszenarien (\dq Anwendungsfall\dq ). Es drückt aus, was das System tun soll, verzichtet dabei aber auf Realisierungsdetails wie Datenstrukturen, Algorithmen etc. Außerdem bildet das Diagramm ab, welche Funktionen welcher Nutzer tatsächlich nutzt. 	
\end{notebox}

\subsection{Übersicht}
Der Anwendungsfall ist ein grundlegendes Konzept der \ac{OOA}. Anwendungsfälle repräsentieren \textit{funktionale Anforderungen} der Kunden an ein System und stellen diese auf einem sehr hohen Abstraktionsniveau dar. Das Anwendungsfalldiagramm klärt die folgenden Fragen für einen Sachverhalt:

\begin{tipbox}{Beispiele}
\begin{itemize}
	\item Was wird beschrieben (\textit{System})?
	\item Wer interagiert mit dem System (\textit{Akteur})?
	\item Was machen die Akteure (\textit{Anwendungsfälle})?
\end{itemize}	
\end{tipbox}

Das Anwendungsfalldiagramm besteht im Vergleich zu den anderen UML-Diagrammen nur aus sehr wenigen Symbolen und wird häufig unterschätzt. Im Folgenden soll uns das Anwendungsfalldiagramm aus Abb. \ref{fig:uc_schuelerverwaltung.drawio} dienen. 

\begin{figure}[h]
  \centering
  \includegraphics[width=0.7\textwidth]{\imgpath/uc_schuelerverwaltung.drawio.png}
  \caption{Anwendungsfalldiagramm Schülerverwaltung}
  \label{fig:uc_schuelerverwaltung.drawio}
\end{figure}

Ein \textit{Anwendungsfall} wird durch eine Ellipse, in der der konkrete Anwendungsfall steht, dargestellt. Er beschreibt das Verhalten, das von dem zu entwickelnden System erwartet wird. Es umfasst eine Menge von Aktionen, die im Kontext der Verwendung durchgeführt werden. Ein Anwendungsfall bringt einen \textit{sichtbaren Nutzen} für einen oder mehrere \textit{Akteure}.  \\ 

Die Menge aller Anwendungsfälle beschreibt die Funktionalität des zu beschreibenden Softwaresystems. Die Anwendungsfälle werden durch ein Rechteck zusammengefasst, welche die Grenzen des zu beschreibenden Systems (hier Schülerverwaltung) aufzeigt. Die hier gezeigten Anwendungsfälle dürfen von verschiedenen \textit{Akteuren} ausgeführt werden. \\ 

Ein \textit{Akteur} ist menschlich (hier durch ein Strichmännchen dargestellt) oder nichtmenschlich (z. B. ein E-Mail-Server, der durch einen Rechner dargestellt wird). Nichtmenschliche Akteure dürfen aber ebenfalls durch Strichmännchen dargestellt werden. \\ 

An der Durchführung des Anwendungsfalls \textbf{Noten eintragen} ist ausschließlich ein Akteur (Klassenlehrer) beteiligt. An dem Anwendungsfall \textbf{Zeugnis ausstellen} \textit{kann} neben dem Klassenlehrer auch ein Sekretär beteiligt sein. \textit{Kann}, nicht \textit{muss} wegen der Multiplizität 0..1. Bei einer Multiplizität von 1 oder höher (bei 1 kann man die Multiplizität weglassen) müssen beide Akteure an dem Anwendungsfall beteiligt sien (z. B. Anwendungsfall Mündliche Prüfung, benötigt Lehrer und Schüler). \\ 

Die Pfeile mit nicht ausgefüllten Pfeilspitzen deuten eine \textit{Generalisierung} an. Im Beispiel sind Klassenlehrer und Sekretär beide Mitarbeiter. Sie erben den Anwendungsfall \textbf{Schülerdaten abfragen} von Mitarbeiter. Zusätzlich können die jeweiligen Akteure (hier sind natürlich Gruppen gemeint) jeweils für ihr Aufgabengebiet weitere Anwendungsfälle bedienen. 

\subsection{Generalisierung}
Die Generalisierung aus Abb. \ref{fig:uc_schuelerverwaltung.drawio} soll noch einmal behandelt werden. In Abb. \ref{fig:uc_schuelerverwaltung2.drawio} ist links dargestellt, dass für den Anwendungsfall \textbf{Schülerdaten abfragen} sowohl der Klassenlehrer als auch der Sekretär nötig sind. Rechts hingegen ist die Generalisierung gezeigt, das heißt: \textbf{Schülerdaten abfragen} dürfen sowohl der Klassenlehrer als auch der Sekretär. Sie müssen aber \textbf{nicht} beide in Aktion treten.   

\begin{figure}[h]
  \centering
  \includegraphics[width=1\textwidth]{\imgpath/uc_schuelerverwaltung2.drawio.png}
  \caption{Anwendungsfalldiagramm Generalisierung vs. mehrfacher Beteiligung}
  \label{fig:uc_schuelerverwaltung2.drawio}
\end{figure}

\subsection{Assoziationen}
Unter einer Assoziation versteht man die Kante, die einen Akteur mit einem Anwendungsfall verbindet. Jeder Akteur muss mit mindestens einem Anwendungsfall verbunden werden, genauso muss jeder Anwendungsfall mit mindestens einem Akteur verbunden sein. Eine Assoziation darf Multiplizitäten aufweisen. Wird eine Multiplizität 0..n (also 0..1, 0..3, 0..n) angegeben, ist die Beteiligung optional, darf dann aber bis zu 1, 3, n Beteiligungen aufweisen. Bei einer Multiplizität von 1..3 muss mindestens ein Akteur teilnehmen, maximal drei. 


\subsection{Beziehungen zwischen Anwendungsfällen}
Wenn ein Anwendungsfall A einen Anwendungsfall B inkludiert, dargestellt als gestrichelter Pfeil mit offener Pfeilspitze von A nach B, der mit dem Schlüsselwort \texttt{<<include>>} gekennzeichnet ist, so wird das Verhalten von B in das Verhalten von A eingebunden. Dabei wird A als Basisanwendungsfall und B als inkludierter Anwendungsfall bezeichnet. Der Basisanwendungsfall benötigt \textbf{immer} das Verhalten des inkludierten Anwendungsfalls, um seine Funktionalität anzubieten. In einer prozeduralen Programmiersprache entspricht die \texttt{<<include>>}-Beziehung dem Aufruf einer Funktion. \\ 

\begin{figure}[h]
  \centering
  \includegraphics[width=0.7\textwidth]{\imgpath/uc_schuelerverwaltung3.drawio.png}
  \caption{Beziehungen zwischen Anwendungsfällen}
  \label{fig:uc_schuelerverwaltung3.drawio}
\end{figure}

In Abb. \ref{fig:uc_schuelerverwaltung3.drawio} inkludiert also der Anwendungsfall \textbf{Lehrveranstaltung anlegen} das \textbf{Lehrer zuweisen} und kann ohne Lehrer nicht angelegt werden. \\ 

Jeder Anwendungsfall bleibt hier einzeln aufrufbar durch den Stundenplaner. Außerdem darf ein Basisanwendugnsfall mit beliebig vielen inkludierte Anwendungsfälle verknüpft sein. Genauso darf ein inkludierter Anwendungsfall mit beliebieg vielen Basisanwendungsfällen verknüpft sein. Hier muss aber darauf geachtet werden, dass kein \textbf{Zyklus} modelliert wird. \\ 

Steht ein Anwendungsfall B mit einem Anwendungsfall A in in einer \texttt{<<extend>>}-Beziehung, darf A das Verhalten von B nutzen, muss aber nicht. B darf also von A aktiviert werden. B bezeichnet man als erweiterter Anwendungsfall.  \\ 

In Abb. \ref{fig:uc_schuelerverwaltung3.drawio} ist das \textbf{Raum buchen} eine optionale Erweiterung von \textbf{Lehrveranstaltung anlegen}, kann aber auch jederzeit einzeln von dem Stundenplaner aufgerufen werden. \\ 

Für jede \texttt{<<extend>>}-Beziehung darf eine Condition formuliert. Eine Condition könnte in dem Anwendungfalldiagramm z. B. wie folgt formuliert werden: \texttt{Condition: {Raum frei}}. Zur Anwendung des Symbols wird auf den Belegsatz in Abb. \ref{fig:ihk_beleg_anwendungsfall} verwiesen. \\ 

Im Zusammenhang der Generalisierung können auch Anwendungsfälle generalisiert werden. Durch einen abstrakten Anwendungsfall (\texttt{<<abstract>>} Anwendungsfall) kann man den generalisierten Anwendungsfall selbst auch unausführbar machen. Dadurch kann man den Implementierer zwingen, dass der \textit{spezialisierte Anwendungsfall} ausgeführt werden muss. Dabei werden alle \texttt{<<include>>}- und \texttt{<<extend>>}-Beziehungen des generalisierten Anwendungsfall vererbt. 

\subsection{Belegsatz}

\begin{figure}[H]
  \centering
  \includegraphics[width=0.9\textwidth]{\imgpath/ihk_beleg_anwendungsfall.png}
  \caption{IHK-Belegsatz für UML-Anwendungsfalldiagramme}
  \label{fig:ihk_beleg_anwendungsfall}
\end{figure}
\newpage
\section{Klassendiagramm}
\todo{Klassendiagramm - Kapitel ergänzen}

\subsection{Belegsatz}
\begin{figure}[H]
  \centering
  \includegraphics[width=0.9\textwidth]{\imgpath/ihk_beleg_klasse.png}
  \caption{IHK-Belegsatz für UML-Klassendiagramme}
  \label{fig:ihk_beleg_klasse.png}
\end{figure}
\newpage
\section{Zustandsdiagramm}
\includepdf[pages=1-3]{\filepath/Zustandsautomat.pdf}
\includepdf[pages=5-8]{\filepath/Zustandsautomat.pdf}

\subsection{Belegsatz}

\begin{figure}[h]
  \centering
  \includegraphics[width=0.7\textwidth]{\imgpath/ihk_beleg_zustand.png}
  \caption{IHK-Belegsatz für UML-Zustandsdiagramme}
  \label{fig:ihk_beleg_zustand}
\end{figure}
\newpage

\section{Aktivitätsdiagramm}

\begin{tipbox}{Beispiele}
	Das Aktivitätsdiagramm dient der Beschreibung einer \textit{Aktivität}. Das kann z. B. sein
	\begin{itemize}
		\item Ein Algorithmus (z. B. Login-Vorgang in einer Anwendung. Dabei kann der Kontrollfluss und Bedingungen klar dargestellt werden)
		\item Ein Geschäftsprozess (z. B. Bestellprozess in einem Onlineshop. Anwendung findet diese Variante in der Softwareanalyse, denn hier werden Entscheidungen und alternative Abläufe mit verschiedenen Rollen berücksichtigt)
		\item parallele Prozesse (z. B. Bestellbearbeitung in einem Shop-System mit parallelen Zweigen (Zahlung prüfen, Rechnung erzeugen, Lagerbestand aktualisieren, Versand vorbereiten))
	\end{itemize}
\end{tipbox}

\subsection{Übersicht}
\todo{Übersicht Aktivitätsdiagramme}
Beschreibung des Aktivitätsdiagramms...
Das folgende Aktivitätsdiagramm zeigt einen einfachen Prozess, um ein verlorenes Passwort wiederzuerlangen.

\begin{figure}[H]
  \centering
  \includegraphics[width=0.6\textwidth]{\imgpath/act_passwortwiderherstellung.drawio.png}
  \caption{Aktivität Passwortwiederherstellung}
  \label{fig:prozess_passwort_vergessen}
\end{figure}

\subsection{Spezielle Sprachenelemente}
\todo{Sprachenelemente beschreiben ..}

\subsection{Belegsatz}
\begin{figure}[H]
	\centering
	\includegraphics[width=1\textwidth]{\imgpath/ihk_beleg_aktivität.png}
	\caption{IHK-Belegsatz für UML-Aktivitätsdiagramme}
	\label{fig:uml_ihk_belegsatz}
\end{figure} 
\newpage
\subsection{Übung: Aktivitätsdiagramme}

% ------------------------------------------------------------


\begin{situationbox}{Situation \cite{Hardy2020}}
	Die Abteilung \dq Desktop-Entwicklung\dq~der IT-Firma ProConsult führt regelmäßig Schulungen für Kunden durch. Die Planung dieser Schulungen wird von der \textit{Vertriebsabteilung} durchgeführt. Für den Ablauf wurde vor einiger Zeit ein Programmablaufplan erstellt.
\end{situationbox}

 \begin{aufgabe}\label{task:uml_beginner}
 		\underline{Analysiere} den dargestellten Programmablaufplan und \underline{erstelle} ein \textbf{UML-Aktivitäts-diagramm}. 
\end{aufgabe}
 		
\begin{taskbox}{ Weitere Infos zur Aufgabenlösung}
	\begin{itemize}
		\item Das Buchen der Dozenten und Räume soll parallel dargestellt werden
		\item Die Aktionen sollen in entsprechenden Swimlanes (Verantwortungsbereichen) dargestellt werden
	\end{itemize}
\end{taskbox}
	 
	\begin{figure}[H]
		\centering
		\includegraphics[width=0.4\textwidth]{\imgpath/uml_activity_einfach.jpeg}
		\caption{PAP des Prozesses Schulung für Kunden durchführen \cite{Hardy2020}}
		\label{fig:uml_activity_einfach}
	\end{figure}


\newpage

% ------------------------------------------------------------

\begin{situationbox}{Situation}
Die AMAG Soft AG ist ein Hersteller von Batterien und Akkus für E-Bikes, E-Scooter, Pedelecs und E-Roller. Zur jährlichen Aktionärsversammlung treffen sich die Aktionäre am Stammsitz. Die AMAG Soft AG möchte ihren Aktionären zukünftig einen E-Service mit verschiedenen Anwendungsmöglichkeiten anbieten. Zur Realisierung dieses Projekts \dq Aufbau einer E-Service-Plattform\dq~für Aktionäre der AMAG Soft AG wird eine Projektgruppe gegründet, in der Sie mitarbeiten sollen. 
\end{situationbox}

 \begin{aufgabe}\label{task:uml_mittel}
 		Der Einladungsvorgang zur Aktionärsversammlung soll mithilfe eines UML-Aktivitäts-diagramms dargestellt werden. \underline{Ergänze} das vorbereitete UML-Diagramm. \footnote{Quelle: IHK-Prüfung}
 \end{aufgabe}
 
 \begin{figure}[H]
	\centering
	\includegraphics[width=1\textwidth]{\imgpath/uml_activity_einfach_ihk_2024.png}
%	\caption{UML}
%	\label{fig:uml_activity_einfach_ihk_2024}
\end{figure}
 

\begin{taskbox}{Informationen zum Prozess}
	\begin{itemize}
		\item Ein Mitarbeiter des Aktionärssekretariats versendet die Einladungen gleichzeitig per Post und über die E-Service-Plattform.
		\item Wenn ein Aktionär einen Zugang zur E-Service-Plattform hat, meldet er sich dort an und kann dabei zwischen virtueller Teilnahme und Präsenzteilnahme wählen. Hat der Aktionär keinen E-Service-Zugang, kann er nur in Präsenz teilnehmen und die Anmeldung per Post versenden. Diese postalische Anmeldung wird vom Mitarbeiter erfasst, (Der Fall, dass von einem Aktionär keine Anmeldung erfolgt, soll ignoriert werden.)
		\item Unmittelbar nach jeder Online-Anmeldung zur Präsenzteilnahme bzw. Erfassung durch den Mitarbeiter im Aktionärssekretariat wird durch das System ein Platz für den Aktionär reserviert.
		\item Bei virtueller Teilnahme wird kein Platz reserviert.
		\item Nach der Reservierung versendet das System entweder die Platzkarten zusammen mit der Tagesordnung an den Aktionär oder bei virtueller Teilnahme die Zugangsdaten und die Tagesordnung.	
	\end{itemize}
\end{taskbox}

% ------------------------------------------------------------
\newpage 
\begin{situationbox}{Situation}
	Die Parkanlage \dq Wilder Garten\dq~in Köln konnte u. a. mit ihrer reichhaltigen Pflanzenvielfalt, [...] und ihren beliebten Attraktionen in der letzten Saison ihre Anziehungskraft als Naherholungsfläche für viele Besucher steigern. Besucher aus allen Altersgruppen und auch Reiseveranstalter zählen zu den Kunden der Parkanlage. 
\end{situationbox}

 \begin{aufgabe}\label{task:uml_schwer}
 		\underline{Analysiere} die untenstehenden Informationen und \underline{erstelle} ein UML-Aktivitätsdiagramm zum Prozess \dq Anlage neuer Pflanzungen\dq \footnote{Quelle: IHK-Prüfung}. Nutze die vorgegebenen Swimlanes.
 \end{aufgabe}
 
\begin{figure}[H]
	\centering
	\includegraphics[width=0.9\textwidth]{\imgpath/uml_activity_schwer_ihk_2023.png}
%	\caption{UML}
%	\label{fig:uml_activity_einfach_ihk_2024}
\end{figure}

\begin{taskbox}{Informationen zum Prozess}
	\textbf{Sie beginnen mit einem Mitarbeiter der Gartenplanung}
	
	\begin{itemize}
		\item Dieser teilt lhnen mit, dass er den Gärtnern den Auftrag zu einer neuen Pflanzung erteilt.
		\item Wenn die Gärtner fertig gepflanzt haben, fertigt er Bilder des neuen Bereichs an. Wenn die Bilder gemacht sind und der Einkauf die Rechnung bezahlt hat, ist der Vorgang abgeschlossen.
	\end{itemize}
	
	\textbf{Informationen eines Gärtners:}
	\begin{itemize}
		\item Wenn der Gärtner einen neuen Pflanzauftrag vom Gartenplaner erhält, stellt er zuerst die Anforderungen zusammen und sagt dem Einkauf welche Pflanzen zu bestellen sind.
		\item Wenn Pflanzen ankommen, bekommt er diese vom Einkauf und bringt die Pflanzen in die Quarantänestation. Der Gärtner prüft täglich die Pflanzen und kontrolliert das Quarantäneende.
		\item Ist eine Pflanze in Quarantäne erkrankt, so muss diese behandelt und die Quarantäne für alle Pflanzen verlängert werden,
		\item sind alle Pflanzen gesund und die Quarantäne abgelaufen pflanzt der Gärtner die Pflanzen.
		\item Nach Abschluss der Pflanzarbeiten macht die Gartenplanung noch Bilder.
	\end{itemize}
	
	\textbf{Der Einkauf gibt Ihnen folgende Hinweise:}
	\begin{itemize}
		\item Der Gärtner teilt uns mit was er braucht, dann bestellen wir diese Pflanzen bei unserem Lieferanten. Bei Anlieferung übergeben wir die Pflanzen an den Gärtner und bezahlen die Rechnung.
		\item Weiterhin beschreibt der Einkauf den Vorgang beim Lieferanten so, dass der Lieferant den Auftrag annimmt, danach die Lieferung zusammenstellt und dann die Pflanzen ausliefert.
	\end{itemize}
\end{taskbox}

\newpage
\section{Sequenzdiagramm}

\begin{notebox}{Einsatzzweck}
Sequenzdiagramme setzt man ein, um das \textit{Inter}objektverhalten, also die Interaktionen zwischen den betiligten Objekten zu beschreiben. 	
\end{notebox}

\begin{tipbox}{Beispiele}
Beispiele für den Einsatz eines Sequenzdiagramms sind:
	\begin{itemize}
		\item Interaktion zwischen Personen (z. B. Lehrer mit Schülern)
		\item Ablauf von Kommunikationsprotokollen (z. B. HTTP)
		\item Kommunikationsabläufe zwischen Anwender und Software
		\item Kommunikationsabläufe bei einem Feueralarm an der Schule
		\item Kommunikaitonsabläufe zwischen Objekten in einem Programm
	\end{itemize}	
\end{tipbox}

\subsection{Übersicht}

Interaktionen beschreiben den Austausch von \textit{Nachrichten} zwischen Interaktionspartnern. Diese Interaktionen beschreiben den Nachrichtenaustausch in einem bestimmten Kontext, also z. B. für die Durchführung eines bestimmten Prozesses. Für die Interkation bestehen dann ggf. bestimmte \textit{Voraussetzungen} und die Interaktion wird durch ein bestimmtes \textit{Ereignis} oder einen vorgegebenen \textit{Zeitpunkt} (z. B. 23 Uhr) ausgelöst. \\

Da ein Interaktionsdiagramm nur in einem bestimmten Kontext definiert ist, werden zumeist nicht alle Interaktionen zwischen den Interaktionspartnern abgedeckt. Sprich: Falls eine bestimmte Interaktion im Diagramm nicht angezeigt wird, heißt das nicht zwangsläufig, dass diese Interkation nicht stattfindet. 

\begin{importantbox}{Abgrenzung zum Aktivitätsdiagramm}
Im Interaktionsdiagramm geht es um die Darstellung von Nachrichten zwischen Interaktionspartnern. Zwar kann man die Nachrichten durch die konkreten Daten anreichern, jedoch geht es nicht um die Darstellung dessen, wie die Daten verändert werden. Dazu würde man eher ein Aktivitätsdiagramm verwenden.	
\end{importantbox}

Als Beispiel soll das Sequenzdiagramm in Abb. \ref{fig:sd_uebung_anmelden.drawio} dienen. Das Diagramm zeigt den Anmeldevorgang für Studenten, welche sich für eine Übung zur Vorlesung anmelden möchten. Die an der Kommunikation beteiligten Kommunikationspartner sind durch Rechtecke dargestellt. Jedes Objekt hat eine \texttt{Lebenslinie}, welche vertikal verläuft. \\ 

Sofern nicht anders definiert, treten Nachrichten von oben nach unten zeitlich nacheinander auf. Zunächst sind die Lebenslinien voneinander unabhängig, erst durch die Übertragung einer Nachricht, synchronisieren sich die Lebenslinien und es gelten für alle der gleiche zeitliche Rahmen. \\ 

\textit{Synchrone Nachrichten} sind solche Nachrichten, auf deren Antwort der Sender wartet. Sie werden durch einen Pfeil mit durchgezogener Linie mit einem \underline{ausgefüllten Pfeil} dargestellt. Die \textit{Antwortnachricht} wird durch eine gestrichelte Linie mit \underline{nicht} ausgefülltem Pfeil symbolisiert. \textit{Asynchrone Nachrichten} sind solche Nachrichten, auf deren Antwort der Sender nicht wartet. Stattdessen geht der Sender anderen Aufgaben nach (Stichwort parallele Prozesse). 

\begin{figure}[h]
  \centering
  \includegraphics[width=0.7\textwidth]{\imgpath/sd_uebung_anmelden.drawio.png}
  \caption{Sequenz Übung anmelden \cite{Seidl2012}}
  \label{fig:sd_uebung_anmelden.drawio}
\end{figure}

\subsection{Kombinierte Fragmente}
Um nun \textbf{Kontrollstrukturen} besser abbilden zu können, gibt es sogenannte \textit{kombinierte Fragmente}. Sie werden durch Rechtecke dargestellt, deren Art durch ein Pentagon in der linken oberen Ecke spezifiziert wird. In dem Pentagon wird der \textit{Operator} eingetragen. In diesem Kurs lernen wir die \textbf{alt}- und \textbf{opt}-Operatoren kennen. \\

Wie in Abb. \ref{fig:sd_uebung_anmelden.drawio} zu sehen, wird der \textbf{alt}-Operator genutzt, um alternative Interaktionen zu beschreiben. Auf einen Operator folgen ein oder mehrere \textit{Operanden}. Sie sind durch gestrichelte horizontale Linien voneinander getrennt. Im Falle des \textbf{alt}-Operators werden Überwachungsbedingungen (Guards) zur Beschreibung der verschiedenen Operanden eingesetzt (z. B. \texttt{[status==ok]}). \\ 

Der \textbf{opt}-Operator zeigt eine optionale Interaktion an. Daneben gibt es auch den \textbf{loop}-Operator für die iterative Interkation (z. B. loop(1,3) für mindestens 1 und maximal 3 Interaktionen) und den \textbf{break}-Operator für die Ausnahme-Interaktion (Fehlerfall).  \footnote{UML bietet außerdem noch u.a. Nebenläufigkeits-Operatoren. Hier empfehle ich \cite{Seidl2012}}

\subsection{Interaktionspartner}
Interaktionspartner sind nach dem \texttt{objekt:Klasse}-Schema, wie bei Objektdiagrammen designed. Auf diese Weise ist es möglich, sowohl für ganz konkrete Objekte, aber auch für Klassen die Interaktionen zu modellieren. Mögliche Kombinationen sind \texttt{objekt}, \texttt{:Klasse}, \texttt{objekt:Klasse}. 

\subsection{Ausführungsspezifikation}
Die Ausführungsspezifikation wird durch ein Sendeereignis und ein Empfangsereignis beschrieben. Dies ist immer dann nötig, wenn zwischen dem Senden einer Nachricht und dem Empfang der beim Emfpänger (und zurück analog) Zeit vergeht. Das kann passieren, wenn eine Übertragung über eine langsame Leitung erfolgt oder die Generierung der Antwort Zeit benötigt (z. B. aufwändige Berechnung eines Ergebnis). In dem Fall zeichnet man einen Balken auf die Lebenslinie, um den Zeitverzug zu symbolisieren. In diesem Skript wird auf das Darstellen der Ausführungsspezifikation bei \dq kleinen Wartezeiten\dq~verzichtet. \todo{Bild einfügen, für Sende- und Empfangsereignis}

\subsection{Belegsatz}
Der Belegsatz der IHK zeigt die meisten der hier beschriebenen UML-Sprachenfeatures und darüber hinausgehende. Sequenzdiagramme haben weitere Sprachenfeatures, vrgl. \cite{Seidl2012}. Der Belegsatz ist in Abb. \ref{fig:ihk_beleg_sequenz} gezeigt.

\begin{figure}[h]
  \centering
  \includegraphics[width=0.7\textwidth]{\imgpath/ihk_beleg_sequenz.png}
  \caption{IHK-Belegsatz für UML-Sequenzdiagramme}
  \label{fig:ihk_beleg_sequenz}
\end{figure}









\newpage

\newpage
\section{Anhang}
\subsection{Rahmen in UML}
Jedes Diagramm hat eine \textit{contents area}. Optional kann jedes Diagramm einen Frame und eine Überschrift haben, beispielsweise wie in Abb. \ref{fig:anhang_rahmen.drawio} dargestellt.

\begin{figure}[H]
	\centering
	\includegraphics[width=0.6\textwidth]{\imgpath/anhang_rahmen.drawio.png}
	\caption{Rahmen in UML}
	\label{fig:anhang_rahmen.drawio}
\end{figure}

Rahmen werden hauptsächlich dort eingesetzt, bei dem das gezeichente Diagramm grafische Rahmenelemente hat wie z. B. entry/exit Punkte für Zustandsdiagramme. Wenn der Rahmen nicht benötigt wird, kann er weggelassen werden. Die Überschrift hat folgende Syntax:

\begin{verbatim}
[<kind>]<name>[<parameters>]
\end{verbatim}

In den meisten Fällen reicht die Art (kind) und der Name des Diagramms. Die Art muss dazu nicht ausgeschrieben werden, man kann auch Kürzel nutzen, wie sie im Standard definiert sind:
\begin{itemize}
	\item \textbf{act} activity frame
	\item \textbf{cmp} component frame
	\item \textbf{dep} deployment frame
	\item \textbf{sd} interaction frames
	\item \textbf{pkg} package frames
	\item \textbf{stm} state machine frames
	\item \textbf{uc} use case frames
	\item (\textbf{cd} class diagramm frames)
\end{itemize}
\textbf{Hinweis:} Eine Definition für Klassendiagramme \textbf{cd} ist im Standard nicht definiert, in diesem Kurs einigen wir uns darauf, dass die Abkürzung gültig is. Der aktuell gültige Standard ist hier als PDF herunterladbar: \url{https://www.omg.org/spec/UML/}. Die Informationen aus diesem Kapitel sind aus \textit{Annex A: Diagrams} entnommen.

\newpage

\newpage
\bibliographystyle{alpha} % Stil, z. B. plain, alpha, abbrv, unsrt
\bibliography{library}    % Name der .bib-Datei OHNE .bib

\end{document}
